\documentclass[11pt]{article}

\usepackage[T1]{fontenc}
\usepackage[a4paper,margin=30mm]{geometry}
\usepackage{amsmath,amssymb,amsthm,mathtools}
\usepackage{booktabs}
\usepackage{microtype}
\usepackage[hidelinks]{hyperref}

\newtheorem{theorem}{Theorem}[section]
\newtheorem{proposition}[theorem]{Proposition}
\newtheorem{corollary}[theorem]{Corollary}
\newtheorem{lemma}[theorem]{Lemma}
\theoremstyle{remark}
\newtheorem{remark}[theorem]{Remark}

\DeclareMathOperator{\ord}{ord}
\DeclareMathOperator{\rad}{rad}
\newcommand{\N}{\mathbb N}
\newcommand{\eps}{\varepsilon}

\title{Super-Wieferich Depth Layers in Canonical Mersenne Blocks}
\author{ChatGPT}
\date{September 1, 2026}

\begin{document}
\maketitle

\begin{abstract}
We isolate the contribution of base-two super-Wieferich depth to the
powerful part of the Mersenne numbers $2^m-1$.  Inside each exact-order
fibre, the logarithm of the genuinely deep factor has an exact finite
layer-cake expansion.  Truncating this expansion at a moving depth gives a
support term and a high-depth tail, and hence both pointwise and
divisor-average sufficient criteria for the Mersenne power loss to be
$o(m)$.  We also give a complete certificate at the prime $3511$: its
base-two order is the odd number $1755$, and its canonical valuation is
exactly two.  This refutes an even-order strengthening sometimes suggested
by factorizations of $2^d-1$, but it does not refute the super-Wieferich
depth criterion.  Two independent finite enumerations through $10^7$ are
recorded only as computational checks.  The required fixed-base
distribution estimates remain open, so the results below do not prove or
disprove the $abc$ conjecture.
\end{abstract}

\noindent\textbf{Keywords.}
Mersenne numbers; cyclotomic polynomials; Wieferich primes;
multiplicative order; lifting the exponent; $abc$ conjecture.

\section{Exact-order blocks}
\label{sec:blocks}

Let
\[
  M_m=2^m-1,
  \qquad
  W_m=\frac{M_m}{\rad(M_m)}.
\]
For an odd prime $q$, put
\begin{equation}
  d_q=\ord_q(2),
  \qquad
  w_q=v_q(2^{d_q}-1).
  \label{eq:canonical-data}
\end{equation}
Since $d_q\mid q-1$ and $q\nmid(q-1)/d_q$, the odd-prime lifting-the-
exponent formula gives
\begin{equation}
  v_q(2^{q-1}-1)=w_q.
  \label{eq:fermat-depth}
\end{equation}
Thus $w_q\ge2$ is precisely the base-two Wieferich condition, while
$w_q\ge3$ is the base-two super-Wieferich condition.

For each $d\ge1$, define the canonical exact-order block and its two depth
components by
\begin{align}
 E_d&=\prod_{d_q=d}q^{w_q-1},
 &T_d&=\prod_{\substack{d_q=d\\w_q\ge2}}q,
 \label{eq:E-T}\\
 S_d^{(3)}&=\prod_{\substack{d_q=d\\w_q\ge3}}q,
 &D_d&=\prod_{\substack{d_q=d\\w_q\ge3}}q^{w_q-2}.
 \label{eq:S-D}
\end{align}
All these products are finite: every prime in the order-$d$ fibre divides
$2^d-1$.  Comparing the exponent of each prime gives the exact identity
\begin{equation}
                         E_d=T_dD_d.
  \label{eq:E-split}
\end{equation}
The factor $T_d$ records one copy of every repeated exact-order prime;
$D_d$ records only the multiplicity beyond the second copy.

For completeness, the order blocks recover the whole Mersenne power loss
after an elementary lifting factor is retained.  Set
\begin{equation}
  L_m=\prod_{q\mid M_m}q^{v_q(m/d_q)}.
  \label{eq:lifting-factor}
\end{equation}
Here $d_q\mid m$ whenever $q\mid M_m$.

\begin{proposition}[Exact Mersenne ledger]
\label{prop:ledger}
For every $m\ge1$,
\begin{equation}
  W_m=L_m\prod_{d\mid m}E_d,
  \qquad L_m\mid m.
  \label{eq:exact-ledger}
\end{equation}
Consequently $0\le\log L_m\le\log m=o(m)$.
\end{proposition}

\begin{proof}
If $q\mid M_m$, then $d_q\mid m$, and odd-prime LTE gives
\[
 v_q(M_m)=w_q+v_q(m/d_q).
\]
After the radical removes one copy of $q$, the first term contributes
$q^{w_q-1}$ to the unique block $E_{d_q}$, and the second contributes its
factor to $L_m$.  Multiplication over the prime support proves the first
identity in \eqref{eq:exact-ledger}.  Moreover, $d_q\mid q-1$ implies
$q\nmid d_q$, hence $v_q(m/d_q)=v_q(m)$.  The product $L_m$ therefore takes
the full $q$-primary component of $m$ for only a subset of the prime
divisors of $m$, and so $L_m\mid m$.
\end{proof}

Murty and S\'eguin identify square and higher-power divisors of cyclotomic
values with Wieferich and super-Wieferich primes at the unique exact-order
index \cite[Lemmas~4.2--4.3]{MurtySeguin2019}.  Their results do not give a
weighted upper bound within a moving exact-order fibre.  Fellini and Murty
obtain quantitative lower bounds for non-Wieferich prime ideals under a
number-field $abc$ conjecture or under finiteness of the relevant
super-Wieferich primes \cite[Theorems~1.2--1.4]{FelliniMurty2026}.  Neither
hypothesis supplies an unconditional estimate for the fixed-base depth
masses studied below.

\section{The exact depth layer cake}
\label{sec:layer-cake}

For $j\ge3$, define the logarithmic depth layer
\begin{equation}
  \Theta_{d,j}
   =\sum_{\substack{d_q=d\\w_q\ge j}}\log q.
  \label{eq:depth-layer}
\end{equation}

\begin{proposition}[Finite layer-cake identity]
\label{prop:layer-cake}
For every $d$,
\begin{equation}
                    \log D_d=\sum_{j\ge3}\Theta_{d,j}.
  \label{eq:layer-cake}
\end{equation}
The sum is finite.  More precisely, if $M\ge w_q$ for every prime in the
order-$d$ super-Wieferich support, then the right-hand side may be
truncated at $j=M$.
\end{proposition}

\begin{proof}
A fixed prime $q$ contributes $(w_q-2)\log q$ to $\log D_d$.  It occurs
once in every layer with $3\le j\le w_q$, and therefore contributes
\[
              \sum_{j=3}^{w_q}\log q=(w_q-2)\log q.
\]
The prime support is finite, and every valuation is finite.  Interchanging
the resulting two finite sums and then summing the displayed identity over
$q$ proves \eqref{eq:layer-cake}.  The same calculation proves the stated
finite truncation.
\end{proof}

The identity distinguishes the first super-Wieferich layer
$\Theta_{d,3}=\log S_d^{(3)}$ from the accumulated mass at higher depths.
A small first layer alone does not control the full sum unless the
high-depth layers are controlled as well.

\section{Moving-depth truncation}
\label{sec:truncation}

For an integer $K\ge3$, put
\begin{equation}
 R_d(K)=\sum_{\substack{d_q=d\\w_q\ge3}}
             (w_q-K)_+\log q,
 \qquad (x)_+=\max\{x,0\}.
 \label{eq:tail-definition}
\end{equation}

\begin{proposition}[Exact truncation and support bound]
\label{prop:truncation}
For every $K\ge3$,
\begin{align}
 \log D_d
 &=\sum_{\substack{d_q=d\\w_q\ge3}}
      \min\{w_q-2,K-2\}\log q+R_d(K),
 \label{eq:exact-truncation}\\
 \log D_d
 &\le (K-2)\log S_d^{(3)}+R_d(K).
 \label{eq:threshold-bound}
\end{align}
Moreover,
\begin{equation}
                    R_d(K)=\sum_{j\ge K+1}\Theta_{d,j}.
 \label{eq:tail-layers}
\end{equation}
\end{proposition}

\begin{proof}
For integers $w\ge0$ and $K\ge2$, truncated subtraction gives
\begin{equation}
 w-2=\min\{w-2,K-2\}+(w-K)_+
 \label{eq:integer-split}
\end{equation}
whenever $w\ge2$.  Multiplying \eqref{eq:integer-split} by the nonnegative
number $\log q$ and summing over the super-Wieferich support proves
\eqref{eq:exact-truncation}.  The first coefficient is at most $K-2$;
summing that pointwise inequality yields \eqref{eq:threshold-bound}.
Finally, $w-K$ is the number of integers $j$ with $K+1\le j\le w$.
The finite sum interchange used in Proposition~\ref{prop:layer-cake}
therefore proves \eqref{eq:tail-layers}.
\end{proof}

\begin{corollary}[Frequency--depth dichotomy]
\label{cor:dichotomy}
If $K>2$ and $\log D_d\ge A$, then
\begin{equation}
 \frac{A}{2(K-2)}\le\log S_d^{(3)}
 \qquad\text{or}\qquad
 \frac A2\le R_d(K).
 \label{eq:dichotomy}
\end{equation}
\end{corollary}

\begin{proof}
If the first alternative fails, then
$(K-2)\log S_d^{(3)}<A/2$.  Combining this strict inequality with
\eqref{eq:threshold-bound} and $\log D_d\ge A$ forces
$R_d(K)>A/2$.
\end{proof}

Thus a large deep factor cannot be explained merely by the existence of
some super-Wieferich primes.  At every chosen threshold it forces either
large one-copy support mass or large mass above that threshold.

\section{Sufficient asymptotic gates}
\label{sec:gates}

All little-oh statements in this section are taken along positive integers
tending to infinity.

\begin{theorem}[Moving-threshold gate]
\label{thm:moving-threshold}
Let $K_d\ge3$ be an integer-valued threshold.  If
\begin{align}
 (K_d-2)\log S_d^{(3)}&=o(\varphi(d)),
 \label{eq:support-gate}\\
 R_d(K_d)&=o(\varphi(d)),
 \label{eq:tail-gate}
\end{align}
then
\begin{equation}
                         \log D_d=o(\varphi(d)).
 \label{eq:deep-conclusion}
\end{equation}
If, in addition, $\log T_d=o(\varphi(d))$, then
\begin{equation}
 \log\frac{2^m-1}{\rad(2^m-1)}=o(m).
 \label{eq:mersenne-endpoint}
\end{equation}
\end{theorem}

\begin{proof}
Apply \eqref{eq:threshold-bound} with $K=K_d$.  Its two nonnegative
right-hand terms are little-oh of $\varphi(d)$ by
\eqref{eq:support-gate}--\eqref{eq:tail-gate}, proving
\eqref{eq:deep-conclusion}.  The exact identity \eqref{eq:E-split} then
gives $\log E_d=o(\varphi(d))$ under the additional hypothesis.

For clarity, the last implication uses only the standard divisor identity
\[
                         \sum_{d\mid m}\varphi(d)=m.
\]
Given $\eps>0$, choose $D$ so that
$\log E_d\le\eps\varphi(d)$ for $d\ge D$, and put
$C_D=\sum_{d<D}\log E_d$.  Then
\[
 \sum_{d\mid m}\log E_d\le C_D+\eps m.
\]
Proposition~\ref{prop:ledger} and $\log L_m=o(m)$ now imply
\eqref{eq:mersenne-endpoint}.
\end{proof}

The two hypotheses in Theorem~\ref{thm:moving-threshold} are sufficient,
not necessary.  In particular, a uniform depth bound makes
\eqref{eq:tail-gate} vanish for a fixed sufficiently large threshold but
does not by itself prove the support estimate \eqref{eq:support-gate}.
Conversely, sparsity of $S_d^{(3)}$ does not control an unbounded high-depth
tail.

The actual Mersenne endpoint admits a weaker divisor-average interface.

\begin{proposition}[Exact divisor-average endpoint]
\label{prop:divisor-endpoint}
One has
\begin{equation}
 \log W_m=o(m)
 \quad\Longleftrightarrow\quad
 \sum_{d\mid m}\log E_d=o(m).
 \label{eq:divisor-equivalence}
\end{equation}
\end{proposition}

\begin{proof}
Taking logarithms in \eqref{eq:exact-ledger} gives
\[
 \log W_m=\log L_m+\sum_{d\mid m}\log E_d.
\]
All terms are nonnegative and $\log L_m=o(m)$.  Adding or subtracting this
negligible lifting term proves the two implications in
\eqref{eq:divisor-equivalence}.
\end{proof}

Let $s_d$ be a controlled part of the one-copy repeated support and let
$U_d$ be the remaining part, so that
\begin{equation}
                  \log T_d=s_d+U_d,
  \qquad s_d,U_d\ge0.
 \label{eq:small-tail-split}
\end{equation}
In the Mersenne order-block route, $s_d$ is the Brun--Titchmarsh small arm
and $U_d$ collects the surviving one-copy support tails.

\begin{theorem}[Divisor-average threshold gate]
\label{thm:divisor-gate}
Assume
\begin{equation}
                  \sum_{d\mid m}s_d=o(m).
 \label{eq:small-divisor-average}
\end{equation}
Let $K_d\ge3$.  If
\begin{equation}
 \sum_{d\mid m}\left(
   U_d+(K_d-2)\log S_d^{(3)}+R_d(K_d)
 \right)=o(m),
 \label{eq:divisor-gate}
\end{equation}
then \eqref{eq:mersenne-endpoint} holds.
\end{theorem}

\begin{proof}
Equations \eqref{eq:E-split} and \eqref{eq:small-tail-split}, together
with \eqref{eq:threshold-bound}, give the termwise nonnegative bound
\[
 \log E_d
 \le s_d+U_d+(K_d-2)\log S_d^{(3)}+R_d(K_d).
\]
Sum over $d\mid m$ and apply
\eqref{eq:small-divisor-average}--\eqref{eq:divisor-gate}.  The right-hand
side is $o(m)$, so Proposition~\ref{prop:divisor-endpoint} yields
\eqref{eq:mersenne-endpoint}.
\end{proof}

If $s_d=o(\varphi(d))$, then \eqref{eq:small-divisor-average} follows from
$\sum_{d\mid m}\varphi(d)=m$.  Likewise, separate pointwise
$o(\varphi(d))$ estimates for the three terms in
\eqref{eq:divisor-gate} imply that divisor-average condition.  No converse
is asserted: a divisor average can tolerate large values at sparse orders.
Theorem~\ref{thm:divisor-gate} is therefore the weaker interface and is
often the more appropriate target.

\section{An exact odd-order counterexample}
\label{sec:3511}

Factorizations involving $2^{d/2}\pm1$ may suggest restricting repeated
exact-order primes to even $d$.  The second classical base-two Wieferich
prime rules out that restriction.

\begin{proposition}[The prime $3511$]
\label{prop:3511}
The integer $3511$ is prime and
\begin{equation}
   \ord_{3511}(2)=1755,
   \qquad
   v_{3511}(2^{1755}-1)=2.
 \label{eq:3511-conclusion}
\end{equation}
In particular, the exact order is odd.
\end{proposition}

\begin{proof}
First, $3511-1=3510=2\cdot3^3\cdot5\cdot13$.  The congruences
\[
 7^{3510}\equiv1\pmod{3511}
\]
and
\[
\begin{array}{c|cccc}
 r&2&3&5&13\\ \midrule
 7^{3510/r}\bmod3511&3510&2754&1800&494\\
 \gcd(7^{3510/r}-1,3511)&1&1&1&1
\end{array}
\]
give a Pocklington certificate for the primality of $3511$: the known
factor is $F=3510>\sqrt{3511}$.

Direct modular exponentiation gives
\begin{align}
 2^{1755}&\equiv1\pmod {3511^2},
 \label{eq:3511-square}\\
 2^{1755}&\equiv21954602502\not\equiv1\pmod {3511^3},
 \label{eq:3511-cube}\\
 2^{585}&\equiv756\pmod {3511},\\
 2^{351}&\equiv1578\pmod {3511},\\
 2^{135}&\equiv88\pmod {3511}.
 \label{eq:3511-order-checks}
\end{align}
The first congruence shows that the order divides
$1755=3^3\cdot5\cdot13$.  The last three congruences exclude division of
the order by each prime divisor $3$, $5$, and $13$; hence the order is
exactly $1755$.  Finally,
\eqref{eq:3511-square}--\eqref{eq:3511-cube} show that its canonical
valuation is exactly two.
\end{proof}

Proposition~\ref{prop:3511} refutes the precise universal statement
\begin{equation}
 q^2\mid2^{\ord_q(2)}-1
 \quad\Longrightarrow\quad
 2\mid\ord_q(2).
 \label{eq:false-parity-claim}
\end{equation}
The scope is narrow.  Since the valuation at $3511$ is two, this prime is
absent from $S_{1755}^{(3)}$ and contributes nothing to $D_{1755}$.  It is
not a super-Wieferich counterexample and does not contradict any condition
in Theorem~\ref{thm:moving-threshold} or
Theorem~\ref{thm:divisor-gate}.

\section{Finite computational audit}
\label{sec:scan}

A finite scan enumerated every prime at most $10^7$, tested
$2^{q-1}\equiv1\pmod{q^2}$, computed the exact order for every hit, and
then increased the prime-power modulus until the congruence first failed.
The generator used a full Eratosthenes sieve.  An independent verifier used
a segmented sieve and recomputed every Fermat congruence, order certificate,
and depth certificate.

The two enumerations agree on all $664{,}579$ primes in the interval.  Their
complete list of base-two Wieferich hits is
\begin{center}
\begin{tabular}{@{}rrrr@{}}
\toprule
$q$ & $\ord_q(2)$ & $w_q$ & order parity\\
\midrule
1093 & 364  & 2 & even\\
3511 & 1755 & 2 & odd\\
\bottomrule
\end{tabular}
\end{center}
Thus the interval contains no base-two super-Wieferich hit.

This is strictly a finite audit.  The absence of a depth-three hit below
$10^7$ proves neither finiteness nor density zero and supplies no
asymptotic estimate for $S_d^{(3)}$, $R_d(K)$, or $D_d$.  It is used only
to check the explicit certificate in Section~\ref{sec:3511} and the
implementation of the depth definitions.

\section{Remaining analytic gate}
\label{sec:open}

The exact identities reduce the deep obstruction without estimating it.
A weaker decomposed sufficient target is to choose moving thresholds
$K_d\ge3$ and prove
\begin{equation}
 \sum_{d\mid m}\left(
   U_d+(K_d-2)\log S_d^{(3)}+R_d(K_d)
 \right)=o(m),
 \label{eq:open-gate}
\end{equation}
while retaining the already controlled divisor-average small arm.  A
stronger pointwise route would instead establish
\[
 \log T_d=o(\varphi(d)),\qquad
 (K_d-2)\log S_d^{(3)}=o(\varphi(d)),\qquad
 R_d(K_d)=o(\varphi(d)).
\]
Neither Murty--S\'eguin nor Fellini--Murty supplies these fixed-base,
weighted exact-order estimates unconditionally.  The truncation bound is a
one-way sufficient mechanism, and no necessity claim is made.  The
super-Wieferich depth route therefore remains open.

\begin{thebibliography}{9}
\small

\bibitem{FelliniMurty2026}
N.~Fellini and M.~R.~Murty,
\newblock Wieferich primes in number fields and the conjectures of
Ankeny--Artin--Chowla and Mordell,
\newblock \emph{J. Number Theory} \textbf{285} (2026), 209--229.
\newblock \url{https://doi.org/10.1016/j.jnt.2026.01.002}.

\bibitem{MurtySeguin2019}
M.~R.~Murty and F.~S\'eguin,
\newblock Prime divisors of sparse values of cyclotomic polynomials and
Wieferich primes,
\newblock \emph{J. Number Theory} \textbf{201} (2019), 1--22.
\newblock \url{https://doi.org/10.1016/j.jnt.2019.02.016}.

\end{thebibliography}

\end{document}
